Partisan gerrymandering concentrates geographically similar partisan voters
into safe districts through an opaque, ad hoc process driven by incumbent
protection and partisan advantage.  We invert the usual goal of algorithmic
redistricting---which optimises for compactness and neutrality---and ask:
can we \emph{deliberately} generate safe seats using transparent, reproducible
graph-partitioning?  We modify the bisect recursive bisection algorithm to
use partisan-similarity edge weights: edges connecting census tracts with
similar Democratic vote share receive large weights, causing the METIS
partitioner to prefer keeping similar-leaning tracts together rather than
minimising boundary length.  This produces ``partisan similarity districts''
(PSDs) that maximise within-district partisan homogeneity.  Applying this
method to North Carolina and Texas using 2020 presidential election data, we
generate PSD maps that achieve higher within-district partisan homogeneity
than the enacted gerrymandered plans in both states.\singrun\  The mean
within-district Democratic vote standard deviation drops from 12.1 to 7.8
percentage points in NC and from 14.3 to 9.2 percentage points in TX under
PSD weighting.\singrun\  Efficiency Gap analysis shows that PSDs can be
calibrated to distribute safe seats proportionally between parties, unlike
enacted gerrymanders that systematically advantage the mapmaking party.
These results establish PSDs as a useful \emph{legal baseline}: courts
evaluating partisan gerrymandering claims now have an algorithmic reference
for what intentional safe-seat maximisation looks like, and can compare
enacted maps against this neutral upper bound.
